\begingroup\scriptsize\setlength{\tabcolsep}{1.2pt}
\begin{tabular}{@{}llrrrrrrrr@{}}
\hline
Materiale & Forma & $p$ (mbar) & $R$ ($\mu$m) & $P$ (mW) & $T$ (K) & $t_{\rm spin}$ (s) & $\eta_{\rm mecc}$ & $\eta_{\rm lim}$ & $F_{\rm scat}/F_g$ \\
\hline
vaterite & sfera & $10^{-6}$ & 1{,}05 & 0{,}316 & 479{,}7 & 9{,}94 & $1{,}78\times10^{-5}$ & $7{,}65\times10^{-8}$ & 8{,}25 \\
vaterite & sfera & $10^{-6}$ & 1{,}46 & 0{,}139 & 371{,}9 & 63{,}04 & $2{,}80\times10^{-6}$ & $7{,}51\times10^{-8}$ & 1{,}32 \\
vaterite & sfera & $10^{-6}$ & 1{,}46 & 0{,}316 & 420{,}3 & 27{,}7 & $6{,}38\times10^{-6}$ & $7{,}51\times10^{-8}$ & 3{,}02 \\
vaterite & sfera & $10^{-6}$ & 1{,}46 & 0{,}72 & 494{,}4 & 12{,}17 & $1{,}45\times10^{-5}$ & $7{,}51\times10^{-8}$ & 6{,}86 \\
\hline
vaterite & sfera & $10^{-8}$ & 1{,}05 & 0{,}316 & 479{,}7 & 9{,}94 & $1{,}78\times10^{-3}$ & $7{,}65\times10^{-8}$ & 8{,}25 \\
vaterite & sfera & $10^{-8}$ & 1{,}46 & 0{,}139 & 371{,}9 & 63{,}04 & $2{,}80\times10^{-4}$ & $7{,}51\times10^{-8}$ & 1{,}32 \\
vaterite & sfera & $10^{-8}$ & 1{,}46 & 0{,}316 & 420{,}3 & 27{,}7 & $6{,}38\times10^{-4}$ & $7{,}51\times10^{-8}$ & 3{,}02 \\
vaterite & sfera & $10^{-8}$ & 1{,}46 & 0{,}72 & 494{,}4 & 12{,}17 & $1{,}45\times10^{-3}$ & $7{,}51\times10^{-8}$ & 6{,}86 \\
\hline
\end{tabular}
\endgroup

\par\smallskip\noindent\footnotesize 8 celle formali, 0 rotori fisici ($\eta_{\rm lim}<10^{-7}$ in tutte): sono limitate dall'attrito molecolare, non da una coppia utile. Tutte richiedono una trappola o un supporto. Qui $F_g=\rho Vg$ e $F_{\rm scat}=P/c$. Il $\beta$ di Epstein vale per sfere; le forme non sferiche sarebbero marcate con $\beta$ incerto.
